\chapter{Дельта-функции, квантовые поля и законы сохранения}

Простого квантового поля, введенного ранее, достаточно, чтобы в общих чертах описать рассеяние, рождение и уничтожение, но из него можно извлечь и больше. В частности, оно показывает, как законы сохранения входят в процессы с частицами. Полный курс квантовой теории поля занимает гораздо больше времени, чем несколько лекций; здесь задача более узкая. Мы воспользуемся самым простым полем, чтобы механически увидеть, откуда берется сохранение энергии, почему при рассеянии на закрепленной мишени импульс рассеянной частицы не обязан сохраняться и как фазовая инвариантность кодирует сохранение заряда.

Прежде чем возвращаться к полю, понадобятся две математические вещи: дельта-функция Дирака и действие операторов рождения и уничтожения на бра-векторы.

\section{Дельта-функция Дирака}

Обозначим координату \(y\), пока не решая, будет ли она позже означать положение, волновое число или что-то иное. Начнем с обычной функции в виде бугорка, сосредоточенной вблизи \(y=a\). Будем делать этот бугорок все уже и уже. Если бы его высоту держали постоянной, его площадь стремилась бы к нулю, поэтому по мере уменьшения ширины будем увеличивать высоту, сохраняя произведение высоты на ширину равным единице. Предельный объект — это дельта-функция Дирака \(\delta(y-a)\).

Вне \(y=a\) она обращается в нуль, бесконечно остро локализована в \(y=a\) и имеет единичную площадь:
\begin{equation}
\delta(y-a)=0\quad (y\neq a),
\qquad
\int_{-\infty}^{\infty}dy\,\delta(y-a)=1.
\end{equation}
Существенно здесь сочетание локализации и фиксированной площади. Сужение пика не устраняет его, потому что одновременно растет его высота.

\begin{figure}[h]
\centering
\includegraphics[width=0.68\textwidth]{lecture_04_figure_02.png}
\caption{Пик дельта-функции единичной площади с центром в \(y=a\).}
\lectureframe{00:05:50}
\end{figure}

Ту же картину можно идеализировать в виде вертикального шипа:
\begin{figure}[h]
\centering
\begin{tikzpicture}[scale=1.0]
\draw[->] (-3.1,0) -- (3.3,0) node[right] {$y$};
\draw[very thick] (0,0) -- (0,2.9);
\draw (0,-0.15) node[below] {$y=a$};
\node at (1.35,1.65) {$\delta(y-a)$};
\end{tikzpicture}
\caption{Идеализованная локализация \(\delta(y-a)\) в нуле ее аргумента.}
\end{figure}

Дельта-функция естественно возникает в интегралах Фурье. Вернемся на мгновение к периодическому пространственному интервалу полной длины \(L\). Периодическая мода имеет вид
\begin{equation}
e^{ikx},
\end{equation}
где \(k\) — одно из допустимых волновых чисел
\begin{equation}
k_n=\frac{2\pi n}{L},
\qquad n\in\mathbb Z.
\end{equation}
Интервал можно задать как \(0\leq x<L\), но симметричнее поместить начало координат в середину и использовать
\begin{equation}
-\frac L2\leq x<\frac L2.
\end{equation}
Концы интервала отождествляются: переход через \(L/2\) возвращает к \(-L/2\).

Проинтегрируем одну допустимую моду по всему периодическому интервалу:
\begin{equation}
I_L(k)=\int_{-L/2}^{L/2}dx\,e^{ikx}.
\end{equation}
Если \(k=0\), подынтегральное выражение тождественно равно единице, и ответ равен \(L\). Если \(k\neq0\) — допустимое значение, положительные и отрицательные части осцилляции взаимно компенсируются на целом числе периодов. Следовательно,
\begin{equation}
I_L(k)=
\begin{cases}
L,&k=0,\\[4pt]
0,&k\neq0,
\end{cases}
\qquad
k=\frac{2\pi n}{L}.
\end{equation}

Теперь будем рассматривать \(I_L(k)\) как функцию на дискретном наборе допустимых значений \(k\). Расстояние между соседними значениями равно
\begin{equation}
\Delta k=\frac{2\pi}{L}.
\end{equation}
При \(k=0\) функция имеет высоту \(L\); при любом другом допустимом значении она равна нулю. Чтобы изобразить этот дискретный пик, придадим ему ширину, равную одному шагу по \(k\). Тогда его площадь равна
\begin{equation}
(\text{height})(\text{width})
=L\frac{2\pi}{L}=2\pi.
\end{equation}
По мере роста \(L\) высота растет, шаг уменьшается, а площадь остается равной \(2\pi\). В пределе бесконечного объема пик превращается в \(2\pi\delta(k)\), что дает формальное тождество Фурье
\begin{equation}
\int_{-\infty}^{\infty}dx\,e^{ikx}=2\pi\delta(k).
\end{equation}
Переменные \(x\) и \(k\) входят в экспоненту симметрично. Если поменять их ролями, получится сопутствующее тождество
\begin{equation}
\int_{-\infty}^{\infty}dk\,e^{ikx}=2\pi\delta(x).
\end{equation}
Эти два элементарных правила скоро превратят интеграл по всем временам взаимодействия в условие сохранения энергии.

\section{Бра, кеты и операторы, действующие влево}

Кет — это символическое обозначение квантового состояния. Для нашей нынешней цели начальное состояние будем записывать как кет, а конечное — как бра. Соединяя бра и кет, получаем скалярное произведение, то есть число. Именно эти числа или их модули в квадрате в конечном счете и можно сравнивать с экспериментом.

\begin{classroomqa}
\lecturetimestamp{00:15:30}
\audiencequestion{А слово ``кет'' вообще уже вводилось раньше?}
\lecturerresponse{Да. Кет — это одна половина бра-кет-нотации; другая половина, которая вводится сейчас, — это бра.}
\end{classroomqa}

Для гармонического осциллятора состояния с заданным числом квантов можно записывать либо как кеты \(\lvert n\rangle\), либо как бра \(\langle n\rvert\). Их скалярное произведение равно
\begin{equation}
\langle m\vert n\rangle=\delta_{mn}
=\begin{cases}
1,&m=n,\\
0,&m\neq n.
\end{cases}
\end{equation}
Дельта Кронекера \(\delta_{mn}\) — дискретный аналог дельта-функции Дирака.

Операторы рождения и уничтожения действуют на кеты обычным образом:
\begin{align}
a^\dagger\lvert n\rangle&=\sqrt{n+1}\,\lvert n+1\rangle,\\
a\lvert n\rangle&=\sqrt n\,\lvert n-1\rangle.
\end{align}
Операторы рождения и уничтожения могут также обозначаться \(a^+\) и \(a^-\); здесь повсюду будут использоваться \(a^\dagger\) и \(a\).

Как те же операторы должны действовать на бра? В обозначениях оператор пишется справа от бра, когда он действует влево. Определение выбирается так, чтобы матричный элемент имел одно и то же значение независимо от того, действует ли оператор сначала на кет или сначала на бра. Рассмотрим
\begin{equation}
\langle n\rvert a^\dagger\lvert m\rangle.
\end{equation}
При действии на кет получаем
\begin{align}
\langle n\rvert a^\dagger\lvert m\rangle
&=\sqrt{m+1}\,\langle n\vert m+1\rangle\\
&=\sqrt{m+1}\,\delta_{n,m+1}.
\end{align}
Это отлично от нуля только если \(n=m+1\), или, что то же самое, если \(m=n-1\). Чтобы, действуя на бра, воспроизвести то же условие, \(a^\dagger\) должен понижать его индекс:
\begin{equation}
\langle n\rvert a^\dagger
=\sqrt n\,\langle n-1\rvert.
\end{equation}
Числовой множитель равен \(\sqrt n\), поскольку для ненулевого матричного элемента \(n=m+1\). Оператор уничтожения подчиняется дополнительному правилу
\begin{equation}
\langle n\rvert a
=\sqrt{n+1}\,\langle n+1\rvert.
\end{equation}
Таким образом, \(a^\dagger\) повышает кет, но понижает бра, тогда как \(a\) понижает кет, но повышает бра.

Преимущество видно на матричном элементе оператора числа
\begin{equation}
\langle n\rvert a^\dagger a\lvert n\rangle.
\end{equation}
При действии вправо
\begin{align}
a\lvert n\rangle&=\sqrt n\,\lvert n-1\rangle,\\
a^\dagger\lvert n-1\rangle&=\sqrt n\,\lvert n\rangle,
\end{align}
поэтому
\begin{equation}
\langle n\rvert a^\dagger a\lvert n\rangle
=n\langle n\vert n\rangle=n.
\end{equation}
Если же действовать влево,
\begin{align}
\langle n\rvert a^\dagger&=\sqrt n\,\langle n-1\rvert,\\
\langle n-1\rvert a&=\sqrt n\,\langle n\rvert,
\end{align}
и снова получается тот же ответ. Не нужно решать, какое из направлений более фундаментально; исчисление устроено так, что в любом случае получается одно и то же число.

\section{Квантовое поле в пространстве и времени}

Причина всей этой системы обозначений — физика элементарных частиц. На микроскопических масштабах главный экспериментальный прием состоит в том, чтобы сталкивать частицы и наблюдать, что получается. Начальное состояние с одной или двумя частицами может перейти в состояние с 4, 5, 7, 9 или каким-то другим числом частиц. Квантовые поля, построенные из операторов рождения и уничтожения, дают язык для описания таких процессов.

Сначала будем работать в одном пространственном измерении. Ничто существенное не мешает использовать три измерения: \(x\) и \(k\) просто становятся трехкомпонентными векторами. Удобно также пользоваться единицами, в которых
\begin{equation}
\hbar=1.
\end{equation}
Тогда импульс и волновой вектор имеют одно и то же численное значение,
\begin{equation}
\vec p=\vec k.
\end{equation}

\begin{figure}[h]
\centering
\includegraphics[width=0.58\textwidth]{lecture_04_figure_03.png}
\caption{Естественные единицы \(\hbar=1\) и отождествление \(\vec p=\vec k\).}
\lectureframe{00:36:56}
\end{figure}

Примем соглашение, что \(\psi\) уничтожает, а \(\psi^\dagger\) рождает:
\begin{align}
\psi(x)&=\sum_k a(k)e^{ikx},\\
\psi^\dagger(x)&=\sum_k a^\dagger(k)e^{-ikx}.
\end{align}
Суммирование идет по \(k=2\pi n/L\) в периодическом пространстве и переходит к непрерывному пределу в бесконечном пространстве. Второе уравнение является эрмитово-сопряженным к первому. Его можно также рассматривать как операторный аналог комплексного сопряжения.\footnote{На доске разложение поля рождения \(\sum_k a^\dagger(k)e^{-ikx}\) временно обозначается как \(\Psi\), а позднее возникает неуверенность в расстановке dagger. Здесь обозначения согласованно зафиксированы так, что \(\psi\) уничтожает, а \(\psi^\dagger\) рождает, в точности как в последующем расчете рассеяния.}

Почему именно такое определение? На этом этапе лучший ответ таков: оно полезно. Это простая сумма по модам импульса, а ее коэффициенты — как раз операторы, которые убирают или добавляют частицы. Подробное обоснование приходит из динамики и из тех вычислений, которые оно делает возможными.

\editorialnote{Короткий фрагмент источника сразу после введения сопряженного пространственного поля недостаточно разборчив, чтобы на его основе делать дополнительное утверждение о нормировке или интерпретации. Надежное содержание вновь начинается с полезности этого определения.}

Пока что поля зависят от положения, но не от времени. Поля осциллируют, и каждому волновому числу \(k\) соответствует частота \(\omega(k)\). Объект с определенной частотой \(\omega\) имеет временную зависимость
\begin{equation}
e^{i\omega t}=\cos(\omega t)+i\sin(\omega t).
\end{equation}
Присоединим сопряженные фазовые множители к двум разложениям поля:
\begin{align}
\psi(x,t)&=\sum_k a(k)e^{ikx}e^{-i\omega(k)t},\\
\psi^\dagger(x,t)&=\sum_k a^\dagger(k)e^{-ikx}e^{i\omega(k)t}.
\end{align}
Теперь каждая пространственная мода осциллирует с частотой, соответствующей ее волновому числу, и поле зависит и от пространства, и от времени.

\section{Нерелятивистское уравнение поля}

Соотношение между \(\omega\) и \(k\) определяет дифференциальное уравнение, которому удовлетворяет поле. Для одного кванта соотношения волна–частица имеют вид
\begin{equation}
E=\hbar\omega,
\qquad
p=\hbar k.
\end{equation}

\begin{figure}[h]
\centering
\includegraphics[width=0.62\textwidth]{lecture_04_figure_04.png}
\caption{Соотношения энергия–частота и импульс–волновое число \(E=\hbar\omega\) и \(p=\hbar k\).}
\lectureframe{00:41:28}
\end{figure}

При \(\hbar=1\) энергия есть частота, а импульс — волновое число:
\begin{equation}
E=\omega,
\qquad
p=k.
\end{equation}
Для нерелятивистской частицы, движущейся намного медленнее света, кинетическая энергия и импульс связаны соотношением
\begin{equation}
E=\frac{p^2}{2m}.
\end{equation}
Это то же соотношение, которое получается из \(E=\tfrac12 mv^2\) и \(p=mv\). Комбинируя эти соотношения, получаем нерелятивистское дисперсионное соотношение
\begin{equation}
\omega(k)=\frac{k^2}{2m}.
\end{equation}

Для поля рождения разложение по модам имеет вид
\begin{equation}
\psi^\dagger(x,t)=\sum_k a^\dagger(k)e^{-ikx}e^{i\omega(k)t}.
\end{equation}

\begin{figure}[h]
\centering
\includegraphics[width=0.66\textwidth]{lecture_04_figure_05.png}
\caption{Разложение по модам с временной зависимостью, используемое для вывода нерелятивистского уравнения поля.}
\lectureframe{00:43:23}
\end{figure}

Продифференцируем по времени:
\begin{equation}
\dot\psi^\dagger(x,t)
=\sum_k i\omega(k)a^\dagger(k)e^{-ikx}e^{i\omega(k)t}.
\end{equation}
Возьмём одну производную по координате:
\begin{equation}
\frac{\partial\psi^\dagger}{\partial x}
=\sum_k(-ik)a^\dagger(k)e^{-ikx}e^{i\omega(k)t}.
\end{equation}
Между этими производными нет прямой связи первого порядка, потому что производная по времени даёт множитель \(\omega(k)\), тогда как одна пространственная производная даёт множитель \(k\). Чтобы явно увидеть различие, рассмотрим линейное дисперсионное соотношение
\begin{equation}
\omega(k)=k.
\end{equation}
Тогда производная по времени дала бы множитель \(i\omega=ik\), тогда как первая пространственная производная даёт множитель \(-ik\). Для каждой моды
\begin{equation}
\frac{\partial\psi^\dagger}{\partial t}
=-\frac{\partial\psi^\dagger}{\partial x}.
\end{equation}
Такое гипотетическое поле удовлетворяло бы уравнению первого порядка. Здесь, однако, \(\omega=k^2/(2m)\), поэтому нужна вторая производная по координате:
\begin{align}
\frac{\partial^2\psi^\dagger}{\partial x^2}
&=\sum_k(-k^2)a^\dagger(k)e^{-ikx}e^{i\omega(k)t}\\
&=-2m\sum_k\omega(k)a^\dagger(k)e^{-ikx}e^{i\omega(k)t}.
\end{align}

\begin{classroomqa}
\lecturetimestamp{00:47:14}
\audiencequestion{Должен ли здесь действительно стоять знак минус, возникающий из пространственных производных?}
\lecturerresponse{Да. Каждая производная по \(x\) от \(e^{-ikx}\) даёт множитель \(-ik\), а две производные дают \((-ik)^2=-k^2\). Поэтому подстановка \(k^2=2m\omega\) даёт \(-2m\omega\).}
\end{classroomqa}

Если умножить производную по времени на \(i\), получается тот же коэффициент при каждой моде, что и при делении второй пространственной производной на \(2m\):
\begin{align}
i\dot\psi^\dagger
&=-\sum_k\omega(k)a^\dagger(k)e^{-ikx}e^{i\omega(k)t},\\
\frac{1}{2m}\frac{\partial^2\psi^\dagger}{\partial x^2}
&=-\sum_k\omega(k)a^\dagger(k)e^{-ikx}e^{i\omega(k)t}.
\end{align}
Следовательно,
\begin{equation}
i\frac{\partial\psi^\dagger}{\partial t}
=\frac{1}{2m}\frac{\partial^2\psi^\dagger}{\partial x^2},
\end{equation}
или, что эквивалентно,
\begin{equation}
-i\frac{\partial\psi^\dagger}{\partial t}
=-\frac{1}{2m}\frac{\partial^2\psi^\dagger}{\partial x^2}.
\end{equation}

\begin{figure}[h]
\centering
\includegraphics[width=0.72\textwidth]{lecture_04_figure_06.png}
\caption{Производные по времени и координате модового разложения порождающего поля.}
\lectureframe{00:49:06}
\end{figure}

\begin{classroomqa}
\lecturetimestamp{00:49:43}
\audiencequestion{Коэффициент равен \(1/(2m)^2\)?}
\lecturerresponse{Нет. Он равен \(1/(2m)\).}
\end{classroomqa}

Сопряжённое поле уничтожения удовлетворяет знакомому уравнению вида Schr\"odinger
\begin{equation}
i\frac{\partial\psi}{\partial t}
=-\frac{1}{2m}\frac{\partial^2\psi}{\partial x^2}.
\end{equation}
Эти два полевых уравнения эрмитово-сопряжены.\footnote{При проверке знака на доске пространственный член на короткое время был записан с противоположным знаком. Непосредственное дифференцирование выписанных фаз мод даёт сопряжённую пару уравнений, записанную здесь.}

Это уравнение имеет тот же вид, что и уравнение Schr\"odinger из обычной квантовой механики, но объект, который ему удовлетворяет, здесь другой. В обычной квантовой механике волновая функция — это комплекснозначная функция, квадрат модуля которой интерпретируется вероятностно. Здесь \(\psi\) и \(\psi^\dagger\) — операторы: \(\psi\) уничтожает частицу, а \(\psi^\dagger\) создаёт её. Эта сопряжённая пара даёт прямую связь между частицами и полями.

Когда соответствующие моды заняты многими квантами, поле может вести себя приближённо классически, подобно тому как сильно возбуждённый гармонический осциллятор может выглядеть классическим, тогда как осциллятор с одним квантом — нет. Большое число заполнения делает возможным большой классический масштаб и может сделать малыми подходящие относительные флуктуации.

\editorialnote{Само по себе большое число заполнения недостаточно для классического поля. Нужно также фазово-когерентное полуклассическое состояние; состояние с определённым числом частиц не имеет определённой фазы и не обязано иметь ненулевое среднее поле.}

\section{Неподвижная мишень и амплитуда рассеяния}

Как используются эти поля? Рассмотрим простейший возможный процесс рассеяния. Частица встречает настолько тяжёлую мишень, что её отдачей будем пренебрегать. Мишень навсегда закреплена в пространстве. Частица может отразиться, пройти через неё или изменить направление; в трёх измерениях направление вылета может отличаться от направления налёта.

Импульс частицы до и после столкновения не обязан быть одинаковым. Для полной физической системы импульс сохраняется, но недостающий импульс поглощает мишень. Как только движение мишени исключается из математического описания, импульс одной лишь частицы уже не сохраняется.

В идеализированной модели её энергия ведёт себя иначе. Представьте, что вы бросаете теннисный мяч в стену. Реальный теннисный мяч нагревается и деформируется, но если этими эффектами пренебречь, он возвращается с той же кинетической энергией, хотя его импульс меняется на противоположный. Модель неподвижной мишени специально устроена так, чтобы воспроизвести это различие: энергия частицы сохраняется, тогда как её импульс может изменяться.

Изобразим пространство по горизонтали, а время по вертикали. Мишень — это вертикальная прямая \(x=0\). Входящее одночастичное состояние с импульсом \(k_i\) достигает мишени, а исходящее одночастичное состояние с импульсом \(k_f\) покидает её.

\begin{figure}[h]
\centering
\begin{tikzpicture}[scale=1.0]
\draw[->] (-3.2,0) -- (3.2,0) node[right] {space};
\draw[->] (0,-0.3) -- (0,4.2) node[above] {time};
\draw[very thick] (0,0) -- (0,3.8);
\node[left] at (-0.15,3.55) {$x=0$};
\draw[thick,->] (-2.2,0.8) -- (0,2.0);
\node[below] at (-1.45,1.05) {$k_i$};
\draw[thick,->] (0,2.0) -- (2.2,3.2);
\node[above] at (1.5,2.9) {$k_f$};
\fill (0,2.0) circle (1.4pt);
\end{tikzpicture}
\caption{Частица с импульсом \(k_i\) рассеивается на мишени, закреплённой в \(x=0\), и выходит с импульсом \(k_f\).}
\end{figure}

Будем говорить на языке поглощения и повторного испускания. Мишень уничтожает налетающую частицу и сразу же создаёт вылетающую, возможно с другим импульсом. Это лишь математическая модель акта рассеяния; реальный атом, например, может поглотить фотон, некоторое время оставаться возбуждённым, а затем испустить его.

\begin{classroomqa}
\lecturetimestamp{01:02:28}
\audiencequestion{Почему частица должна быть поглощена и немедленно испущена снова?}
\lecturerresponse{В природе это не обязательно. Здесь используется упрощающая модель мгновенного поглощения и повторного испускания в одной точке.}
\end{classroomqa}

Начальное состояние содержит одну частицу в моде \(k_i\) и ни одной в любой другой моде. На языке чисел заполнения каждая компонента равна нулю, кроме компоненты, соответствующей \(k_i\), которая равна единице. Полный список достаточно сократить до вида
\begin{equation}
\lvert k_i\rangle=a^\dagger(k_i)\lvert0\rangle.
\end{equation}
Аналогично,
\begin{equation}
\langle k_f\rvert=\langle0\rvert a(k_f).
\end{equation}

Уничтожение на мишени в момент \(t\) представляется оператором \(\psi(0,t)\); немедленное создание в той же точке и в тот же момент представляется оператором \(\psi^\dagger(0,t)\). Нет никакого выделенного момента времени, в который происходит это событие. Оно может произойти всякий раз, когда налетающая частица достигает мишени, причём сила взаимодействия одинакова во все моменты времени. Чтобы выразить отсутствие особого времени, проинтегрируем по всем \(t\):
\begin{equation}
\int_{-\infty}^{\infty}dt\,\psi^\dagger(0,t)\psi(0,t).
\end{equation}
Действуя на начальный кет, этот оператор создаёт состояние. Его перекрытие с конечным бра есть амплитуда указанного перехода. Вероятность получается только после взятия квадрата модуля амплитуды.

Нужно ввести ещё одно число. Пусть \(g\) измеряет силу взаимодействия между частицей и мишенью. В этой расчётной модели запишем схематическую ведущую амплитуду в виде\footnote{Опущены факторы нормировки состояний и стандартный общий множитель теории возмущений первого порядка; в вычислении используются только операторные правила отбора и область поддержки, совместимая с сохранением энергии.}
\begin{equation}
\mathcal A_{i\to f}
=g\left\langle k_f\left|
\int_{-\infty}^{\infty}dt\,
\psi^\dagger(0,t)\psi(0,t)
\right|k_i\right\rangle.
\end{equation}
Коэффициент \(g\) называется константой связи. Для заряженной мишени, рассеивающей фотон, электрический заряд даёт соответствующую меру силы взаимодействия: с ростом заряда возрастает склонность к рассеянию.

У разных взаимодействий константы связи сильно различаются. Мезон, налетающий на протон, испытывает сильное взаимодействие и имеет сравнительно большую силу рассеяния. У фотона, рассеивающегося на заряженном протоне, электромагнитная константа связи меньше. У нейтрино, взаимодействующего с протоном, она ещё меньше.

\begin{classroomqa}
\lecturetimestamp{01:09:09}
\audiencequestion{Применима ли эта простая модель к электрону?}
\lecturerresponse{Лишь как упрощённый суррогат для нескольких видов рассеяния. Пока её не следует понимать как реалистическую теорию электрона или какой-либо одной конкретной мишени.}
\end{classroomqa}

Гравитационное рассеяние гравитона на ядре было бы чрезвычайно слабым. Во всех этих примерах нет одной универсальной \(g\); в физике частиц имеется много констант связи.

И эффективная константа связи не обязана всегда быть независимой от частоты.

\begin{classroomqa}
\lecturetimestamp{01:10:31}
\audiencequestion{Почему \(g\) берут константой, а не функцией от \(\omega\) или \(k\)?}
\lecturerresponse{Она не обязательно всегда постоянна. Если фотон поглощается атомом, атом некоторое время остаётся возбуждённым, а затем фотон испускается, эффективное рассеяние может зависеть от частоты. Константа \(g\) относится к рассматриваемой здесь мгновенной модели.}
\end{classroomqa}

Пока будем считать \(g\) постоянной и вычислим последствия этого простейшего выбора.

\begin{classroomqa}
\lecturetimestamp{01:12:32}
\audiencequestion{Если \(g=1\), должна ли вероятность рассеяния быть равной единице?}
\lecturerresponse{Нет. Вычисляемое выражение — лишь ведущее приближение при малом \(g\). Если константа связи не мала, нужна теория возмущений более высокого порядка; вероятность больше единицы означала бы нарушение одноступенчатого приближения, а не физическую вероятность.}
\end{classroomqa}

\section{Свёртывание сумм по модам}

При \(x=0\) пространственные фазовые множители равны единице, поэтому
\begin{align}
\psi(0,t)&=\sum_k a(k)e^{-i\omega(k)t},\\
\psi^\dagger(0,t)&=\sum_\ell a^\dagger(\ell)e^{i\omega(\ell)t}.
\end{align}

\begin{classroomqa}
\lecturetimestamp{01:13:58}
\audiencequestion{Поскольку мишень находится в \(x=0\), следует ли опустить пространственную экспоненту?}
\lecturerresponse{Да. На мишени \(e^{ikx}=e^{ik0}=1\), и сопряжённый пространственный множитель тоже равен единице.}
\end{classroomqa}

Используйте разные фиктивные индексы \(k\) и \(\ell\) в двух суммах. Подстановка даёт
\begin{align}
\mathcal A_{i\to f}
&=g\int_{-\infty}^{\infty}dt
\sum_{\ell,k}
e^{i\omega(\ell)t}e^{-i\omega(k)t}
\langle k_f\rvert a^\dagger(\ell)a(k)\lvert k_i\rangle.
\end{align}
Какой член в сумме по \(k\) может дать ненулевой вклад? Оператор уничтожения даёт нуль в любой пустой моде. Начальное состояние содержит частицу только в моде \(k_i\), поэтому \(k=k_i\). Какой член в сумме по \(\ell\) может дать ненулевой вклад? Действуя влево, \(a^\dagger(\ell)\) должен убрать частицу, представленную конечным бра, поэтому \(\ell=k_f\). Обе суммы схлопываются:
\begin{align}
\mathcal A_{i\to f}
&=g\int_{-\infty}^{\infty}dt\,
e^{i(\omega_f-\omega_i)t}
\langle k_f\rvert a^\dagger(k_f)a(k_i)\lvert k_i\rangle.
\end{align}
Оператор уничтожения переводит начальный одночастичный кет в вакуум,
\begin{equation}
a(k_i)\lvert k_i\rangle=\lvert0\rangle,
\end{equation}
а оператор рождения, действуя влево, переводит конечное одночастичное бра в вакуумное бра,
\begin{equation}
\langle k_f\rvert a^\dagger(k_f)=\langle0\rvert.
\end{equation}
Поэтому их матричный элемент равен
\begin{equation}
\langle0\vert0\rangle=1.
\end{equation}
Весь операторный учёт сводится к единице; коэффициент связи даёт \(g\), а оставшаяся нетривиальная часть — это интеграл по времени:
\begin{equation}
\mathcal A_{i\to f}
=g\int_{-\infty}^{\infty}dt\,
e^{i(\omega_f-\omega_i)t}.
\end{equation}

\editorialnote{Краткий вопрос из аудитории в 01:18:40 не удаётся восстановить по источнику. В ответе проводится различие между этим одночастичным вычислением и релятивистской или многочастичной проблемой; более точно восстановить обмен репликами надёжно нельзя.}

Теперь ранее полученное тождество для дельта-функции применяется непосредственно:
\begin{equation}
\int_{-\infty}^{\infty}dt\,
e^{i(\omega_f-\omega_i)t}
=2\pi\delta(\omega_f-\omega_i).
\end{equation}
Следовательно,
\begin{equation}
\boxed{
\mathcal A_{i\to f}
=2\pi g\,\delta(\omega_f-\omega_i)
}.
\end{equation}
Амплитуда обращается в нуль, если только не
\begin{equation}
\omega_f=\omega_i.
\end{equation}
Поскольку \(\hbar=1\), это есть \(E_f=E_i\). Сохранение энергии появилось потому, что все моменты взаимодействия рассматривались одинаково. Если бы интеграл по времени был ограничен конечным интервалом, он дал бы размытую функцию, а не точную дельта-функцию. Именно точная инвариантность относительно сдвигов времени даёт точное условие сохранения энергии.

\begin{classroomqa}
\lecturetimestamp{01:21:30}
\audiencequestion{Чтобы получить вероятность, амплитуду нужно возвести в квадрат. Что происходит с дельта-функцией при возведении в квадрат?}
\lecturerresponse{Квадрат дельта-функции требует дополнительной осторожности. Нужный здесь вывод надёжен: результат равен нулю, если начальная и конечная энергии не равны. На носителе условия сохранения энергии схематический префактор после возведения в квадрат содержит \(4\pi^2g^2\).}
\end{classroomqa}

\editorialnote{Нормированную вероятность перехода или скорость перехода нельзя получить, просто отбросив \(\delta(\omega_f-\omega_i)^2\). Для этого нужны регуляризация по конечному времени и согласованная нормировка состояний. Поэтому \(4\pi^2g^2\) сохраняется только как схематический квадратный префактор, подчеркнутый в источнике.}

Ничто в этом взаимодействии с точечной мишенью, кроме энергетического условия, не различает один конечный импульс от другого. Поэтому в трёх измерениях модель приписывает одинаковый ведущий вес всем вылетающим импульсам с одной и той же энергией---эквивалентно, всем направлениям вылета на сфере фиксированной энергии. Эта изотропия — свойство именно этого очень простого точечного рассеивателя, а не любой реальной мишени.

Вычисление дало три связанные идеи. Коэффициент \(g\) измеряет силу взаимодействия. Выражение типа бра-оператор-кет — это амплитуда рассеяния, модуль которой в квадрате входит в вероятность. Самое важное: интегрирование по всем временам порождает энергетическую дельта-функцию; симметрия относительно сдвигов времени обеспечивает сохранение энергии.

\section{Что меняется, когда меняются мишень или коэффициент связи?}

Начало координат \(x=0\) было выбрано лишь для удобства. Поместим мишень в точку \(x=x_{\mathrm T}\). Пространственные множители в \(\psi^\dagger(x_{\mathrm T},t)\psi(x_{\mathrm T},t)\) дают
\begin{equation}
e^{-ik_f x_{\mathrm T}}e^{ik_i x_{\mathrm T}}
=e^{i(k_i-k_f)x_{\mathrm T}}.
\end{equation}
Следовательно, амплитуда принимает вид
\begin{equation}
\mathcal A_{i\to f}(x_{\mathrm T})
=e^{i(k_i-k_f)x_{\mathrm T}}
\mathcal A_{i\to f}(0).
\end{equation}

\begin{classroomqa}
\lecturetimestamp{01:24:37}
\audiencequestion{Можно ли всегда переопределить координату так, чтобы мишень оказалась в \(x=0\)?}
\lecturerresponse{Для этой одной фиксированной мишени — да. Если оставить её в \(x_{\mathrm T}\neq0\), её положение входит в амплитуду как фаза \(e^{i(k_i-k_f)x_{\mathrm T}}\). Умножение на комплексно-сопряжённую фазу устраняет её из вероятности.}
\end{classroomqa}

Действительно,
\begin{equation}
\left|e^{i(k_i-k_f)x_{\mathrm T}}\right|^2=1.
\end{equation}
Амплитуда помнит положение мишени, а эта вероятность для одной мишени — нет.

Теперь обратим эту асимметрию. Предположим, взаимодействие могло бы происходить в любом положении, но только в один выделенный момент времени. Интегрирование по координате дало бы
\begin{equation}
\int_{-\infty}^{\infty}dx\,
e^{i(k_i-k_f)x}
=2\pi\delta(k_i-k_f),
\end{equation}
тогда как фиксация времени убрала бы дельта-функцию энергии.

\begin{classroomqa}
\lecturetimestamp{01:26:41}
\audiencequestion{Если бы было неважно, где происходит взаимодействие, но оно происходило бы только в один фиксированный момент времени, сохранялась бы тогда другая величина?}
\lecturerresponse{Да. Интегрирование по всем положениям обеспечило бы сохранение импульса, тогда как выделенный момент взаимодействия устранил бы точное сохранение энергии. В полной двухчастичной задаче взаимодействие может происходить в любой момент времени и в любом месте.}
\end{classroomqa}

Это также даёт физическую интерпретацию явной зависимости от времени. Если коэффициент связи имеет вид \(g(t)\) или внешнее вмешательство внезапно делает рассеяние возможным, не все моменты времени эквивалентны. Тогда редуцированное описание может обмениваться энергией с той внешней динамикой, которая была опущена.

\begin{classroomqa}
\lecturetimestamp{01:28:29}
\audiencequestion{Константы связи в основном берутся из эксперимента, и что означает зависящий от времени коэффициент связи?}
\lecturerresponse{Да, их значения определяются экспериментально. Явная зависимость от времени обычно означает, что какая-то движущаяся или изменяющаяся часть системы оставлена вне рассмотрения, подобно тому как отдача мишени была оставлена вне модели с фиксированной мишенью. Энергия, которая кажется недостающей, принадлежит этим опущенным степеням свободы.}
\end{classroomqa}

Движущаяся мишень — простейший пример.

\begin{classroomqa}
\lecturetimestamp{01:29:23}
\audiencequestion{Является ли попадание в движущуюся мишень примером, в котором энергия частицы не сохраняется?}
\lecturerresponse{Да. Если предписанное движение мишени не входит в энергетический баланс, она может сообщать частице энергию или отбирать её у неё.}
\end{classroomqa}

Точечный характер мишени — ещё одно упрощение.

\begin{classroomqa}
\lecturetimestamp{01:29:43}
\audiencequestion{Когда размер мишени нужно сравнивать с размером падающей волны или с её длиной волны?}
\lecturerresponse{В данной модели мишень имеет нулевой размер. Конечный размер мишени действительно даёт дополнительные эффекты при сравнении с длиной волны, но это уточнение отложено.}
\end{classroomqa}

Вернёмся теперь к движущейся мишени. Представьте себе стенку, движущуюся навстречу теннисному мячу. Стенка может передать мячу энергию. Если движение стенки задано заранее и её энергия не учитывается, одна только энергия мяча не сохраняется. Математически это зависимость от времени в положении мишени, а не обязательно в \(g\).

Операторы рождения и уничтожения, таким образом, сделали ровно то, что и требовалось: они описывают переход из одного состояния частицы в другое, тогда как пространственно-временная зависимость полей определяет, какие условия сохранения сопровождают этот переход.

\section{Подсчёт заряда и фазовая инвариантность}

Каждому типу частиц соответствует своё поле. Электроны и фотоны, например, не описываются одним и тем же \(\psi\). Чтобы показать подсчёт заряда с помощью уже имеющейся алгебры, будем временно считать, что \(\psi\) обозначает электроноподобное поле: \(\psi\) уничтожает одну такую заряженную частицу, а \(\psi^\dagger\) создаёт одну. Это лишь бозонная замена для такого учёта; подлинное фермионное электронное поле требует алгебры, которая ещё не введена.

Фактор для фиксированной мишени
\begin{equation}
\psi^\dagger\psi
\end{equation}
уничтожает одну частицу и создаёт одну. Один электрон входит и один выходит, так что заряд, переносимый этим типом частиц, не изменяется.

Теперь представьте, что один электрон входит, а два выходят из одной и той же точки. Соответствующее операторное произведение содержало бы одно поле уничтожения и два поля рождения,
\begin{equation}
\psi^\dagger\psi^\dagger\psi.
\end{equation}
Для электронного поля это создало бы один лишний электрон и нарушило бы сохранение электрического заряда. Два электрона на входе и два на выходе можно вместо этого представить как
\begin{equation}
\psi^\dagger\psi^\dagger\psi\psi,
\end{equation}
что содержит одинаковое число полей рождения и уничтожения. Два на входе и три на выходе потребовали бы
\begin{equation}
\psi^\dagger\psi^\dagger\psi^\dagger\psi\psi,
\end{equation}
и опять изменили бы суммарный заряд. Непосредственное правило просто: взаимодействие, сохраняющее заряд, содержит одинаковое число \(\psi\) и \(\psi^\dagger\).

\begin{classroomqa}
\lecturetimestamp{01:34:27}
\audiencequestion{Имеет ли порядок полевых операторов значение для этой проверки сохранения заряда?}
\lecturerresponse{В этой непосредственной проверке важно число множителей \(\psi\) и \(\psi^\dagger\), а не их порядок. Вопрос о том, являются ли произведения с разным порядком в остальном тождественными, — отдельный.}
\end{classroomqa}

Существует математический признак, более глубокий, чем простой подсчёт. Умножим поле на постоянный фазовый множитель:
\begin{equation}
\psi\longrightarrow e^{i\alpha}\psi.
\end{equation}
Поскольку \(\psi^\dagger\) — эрмитово-сопряжённое поле, оно преобразуется с комплексно-сопряжённой фазой,
\begin{equation}
\psi^\dagger\longrightarrow e^{-i\alpha}\psi^\dagger.
\end{equation}
Выражение с одинаковым числом этих двух полей не меняется. Например,
\begin{align}
\psi^\dagger\psi
&\longrightarrow
e^{-i\alpha}e^{i\alpha}\psi^\dagger\psi
=\psi^\dagger\psi,\\
\psi^\dagger\psi^\dagger\psi\psi
&\longrightarrow
e^{-2i\alpha}e^{2i\alpha}
\psi^\dagger\psi^\dagger\psi\psi
=\psi^\dagger\psi^\dagger\psi\psi.
\end{align}
Выражение с неодинаковым числом меняется. Поэтому
\begin{equation}
\psi^\dagger\psi^\dagger\psi
\longrightarrow
e^{-i\alpha}\psi^\dagger\psi^\dagger\psi.
\end{equation}
Разрешённые взаимодействия, сохраняющие заряд, — в точности те, которые инвариантны относительно этого глобального изменения фазы. Фазовая инвариантность и сохранение заряда — два описания одного и того же правила в этом простом языке полей.

Комбинация
\begin{equation}
\psi+\psi^\dagger
\end{equation}
не инвариантна относительно произвольного фазового поворота.

\begin{classroomqa}
\lecturetimestamp{01:37:53}
\audiencequestion{Разве это не вещественная часть \(\psi\)?}
\lecturerresponse{Да. Она пропорциональна вещественной части \(\psi\), а умножение \(\psi\) на фазовый множитель изменяет эту вещественную часть. Следовательно, она не инвариантна относительно фазы.}
\end{classroomqa}

\begin{classroomqa}
\lecturetimestamp{01:38:05}
\audiencequestion{Если электрон ударяет по мишени и испускается много вторичных электронов, как при этом всё ещё может сохраняться заряд?}
\lecturerresponse{Мишень в этом процессе не пассивна. Если атом испускает много электронов, сам атом изменяет свой заряд и состояние. Нужно учитывать его динамику, и саму мишень тоже следует описывать квантовыми полями.}
\end{classroomqa}

На этом этапе видны три соответствия. Инвариантность относительно сдвигов по времени соответствует сохранению энергии. Глобальная фазовая инвариантность соответствует сохранению заряда. Инвариантность относительно пространственных сдвигов соответствует сохранению импульса, хотя полный пространственный расчёт ещё предстоит выполнить.

\section{Порядок операторов и пределы модели низшего порядка}

Для простоты в произведениях выше все поля рождения помещены слева, а все поля уничтожения — справа. В чередующемся произведении числа тех и других всё ещё могут быть равны, и потому оно проходит проверку по подсчёту заряда, но порядок операторов может влиять на выражение и другими способами.

\begin{classroomqa}
\lecturetimestamp{01:40:02}
\audiencequestion{Можно ли тот же тип процесса записать с чередующимися множителями, например \(\psi\psi^\dagger\psi\psi^\dagger\)?}
\lecturerresponse{В нём числа тех и других по-прежнему равны, так что правило для заряда выполняется, но в общем случае это не то же самое выражение, что и выражение, в котором все сопряжённые поля стоят слева. Обсуждение перестановки различных упорядочений откладывается.}
\end{classroomqa}

\editorialnote{Формулировка вопроса из аудитории в 01:41:10 слишком повреждена, чтобы её можно было восстановить. Из сохранившегося ответа повторяется, что равные числа \(\psi\) и \(\psi^\dagger\) означают фазовую инвариантность и сохранение заряда, тогда как неравные числа не означают ни того, ни другого.}

Эти связи между симметриями и законами сохранения являются проявлениями теоремы Нётер. В классической механике эта теорема задаёт общую схему; здесь та же связь была показана непосредственно через фазы и дельта-функции квантовых полей.

Ведущий результат для эксперимента с неподвижной мишенью также имеет очевидную область применимости. Схематическому префактору вероятности, такому как \(4\pi^2g^2\), нельзя доверять, как только он приближается к единице или превосходит её.

\begin{classroomqa}
\lecturetimestamp{01:43:33}
\audiencequestion{Означает ли этот схематический результат, что \(g\) должно быть меньше примерно \(1/(2\pi)\)?}
\lecturerresponse{Либо \(g\) достаточно мало, чтобы этот первый член был полезен, либо модель неполна. Частица может рассеяться один раз, вернуться и рассеяться дважды, затем трижды и так далее. Полная амплитуда представляет собой бесконечный ряд. Только при малом \(g\) простейший член доминирует; в противном случае нужно суммировать много диаграмм Фейнмана.}
\end{classroomqa}

Простые пространственно-временные картинки — это зачатки диаграмм Фейнмана. Их вершины задают элементарные возможности: одной входящей и одной выходящей частице соответствует множитель, такой как \(\psi^\dagger\psi\); двум входящим и двум выходящим — множитель, такой как \(\psi^\dagger\psi^\dagger\psi\psi\). Повторяющиеся вершины порождают более высокие степени константы связи и члены более высокого порядка в ряду.

\begin{classroomqa}
\lecturetimestamp{01:45:39}
\audiencequestion{Может ли константа связи иметь мнимую часть?}
\lecturerresponse{Комплексные константы связи возможны. Неприводимые \(T\)-нечётные фазовые комбинации могут указывать на нарушение инвариантности относительно обращения времени: процесс, обращённый по времени, не обязан иметь амплитуду, связанную этой симметрией.}
\end{classroomqa}

\editorialnote{В этом обмене репликами упоминался параметр КХД \(\theta\), но это вещественный коэффициент, умножающий \(T\)-нечётное взаимодействие, а не мнимая компонента константы связи. Устранимая фаза или абсорбтивная комплексная фаза сама по себе не устанавливает нарушение \(T\). Этот вывод применим только к неприводимым \(T\)-нечётным фазовым комбинациям.}

\begin{classroomqa}
\lecturetimestamp{01:47:02}
\audiencequestion{Может ли выходящая частица испускаться до или после поглощения входящей частицы, а не в тот же самый момент?}
\lecturerresponse{Это справедливый вопрос, но мгновенная модель его не рассматривает. Более общее временное упорядочение отложено.}
\end{classroomqa}

\begin{classroomqa}
\lecturetimestamp{01:47:44}
\audiencequestion{Как сделать эту конструкцию релятивистской?}
\lecturerresponse{Короткое указание состоит в том, чтобы объединить сопряжённые части и включить знаменатель, содержащий квадратный корень из \(\omega(k)\). Полное построение откладывается.}
\end{classroomqa}

\editorialnote{Краткую устную релятивистскую формулу невозможно восстановить достаточно надёжно, чтобы набрать её. Сохранены только перечисленные для неё составные элементы и указание на то, что её рассмотрение отложено.}

Следует также чётко отличать сохранение заряда от сохранения массы покоя.

\begin{classroomqa}
\lecturetimestamp{01:48:22}
\audiencequestion{Сохраняется ли масса покоя в том же структурном смысле, что и электрический заряд?}
\lecturerresponse{Нет. В галилеевой нерелятивистской механике масса является аддитивной сохраняющейся величиной. В релятивистских реакциях сумма масс покоя отдельных частиц не обязана сохраняться. Электрон и позитрон могут аннигилировать в два фотона, тогда как электрический заряд остаётся сохранённым.}
\end{classroomqa}

\editorialnote{Для полной замкнутой релятивистской системы сохраняется полный четырёхимпульс, поэтому инвариантная масса этой системы фиксирована. Пример с аннигиляцией относится к сумме масс покоя отдельных частиц, а не к инвариантной массе замкнутой системы.}

Следовательно, сохранение аддитивной массы — это эффективный нерелятивистский факт, а не фундаментальное правило, стоящее на одном уровне с сохранением электрического заряда. Обнаруженная выше фазовая инвариантность выражает более глубокую симметрию заряда.

Следующие шаги состоят в том, чтобы заменить бозонный суррогат фермионными полями и допустить несколько видов частиц---электроны, мюоны, кварки, фотоны и другие---, каждый со своим собственным полем и своими допустимыми взаимодействиями.

\begin{classroomqa}
\lecturetimestamp{01:50:39}
\audiencequestion{Было ли корректно называть предыдущий пример электронным полем, если до сих пор были введены только бозонные поля?}
\lecturerresponse{Только в смысле бозонной версии или суррогата электрона. Подлинная фермионная теория ещё не была введена.}
\end{classroomqa}